% FILE: 04_implementation_surface.tex
% Project: standards
% Thesis Role: public-standards-compliance-and-gravity-field-mapping

\section{Implementation Surface}
\label{sec:standards-impl}

\subsection{Source Files}
\label{sec:standards-impl-sources}

The 52 source files of the standards project are organized into four functional categories.
All paths are absolute under \texttt{/Users/sac/process-intelligence/standards/}.

\paragraph{Standard placement files (16 files).}
Each placement file is a fully realized schema with no stubs, as confirmed by the
graduation gate. Representative placement files include:
\begin{itemize}
  \item \texttt{xes\_process-intelligence\_placement.md}: XES (IEEE 1849) ontological
        mapping, trace hash chaining formula, chronological invariance constraint, XES-to-OCEL
        and OCEL-to-XES loss policies with signed LossReport schemas.
  \item \texttt{ocel\_process-intelligence\_placement.md}: OCEL 2.0 bipartite graph
        formulation $\mathcal{L}=(\mathcal{E}, \mathcal{O}, \text{E2O}, \text{O2O}, \Delta)$,
        dimensional schema $(\mathcal{OT}, \mathcal{AT}, \lambda)$, SHACL constraint
        references, and multi-object alignment rules.
  \item \texttt{petri\_net\_placement.md}: Places, transitions, arcs, enabling and firing
        rules, 1-boundedness enforcement.
  \item \texttt{wf-net\_placement.md}: Source/sink place requirements, short-circuited
        liveness, soundness criteria grounded in van der Aalst (1998).
  \item \texttt{declare\_placement.md}: 21 LTL$_f$ Declare templates, FSA trace
        acceptance, \texttt{DeclareTemplate} enum variants, violation records.
\end{itemize}

\paragraph{Strategic doctrine files (7 files).}
These files elaborate the competitive and mathematical implications of the standards
registry:
\begin{itemize}
  \item \texttt{PROCESS\_STANDARDS\_GRAVITY\_FIELD.md}: The gravitational model; 12+
        public standards as the center around which all process intelligence capability
        orbits.
  \item \texttt{reverse-lock-in.md}: $W_{\text{std}}$ as Nash equilibrium cage; formal
        competitor behavioral restriction; nullification event ledger schema.
  \item \texttt{reverse-porter-five.md}: Inversion of all five Porter competitive forces
        via Full-Lifecycle Process Intelligence; game-theoretic equilibrium analysis.
  \item \texttt{MAPE\_K\_INTEGRATION.md}: Monitor=OTel ingestion, Analyze=streaming
        conformance, Plan=prescriptive shapes, Execute=\texttt{OcelEvent} actuation.
\end{itemize}

\paragraph{M\&A integration files (4 files).}
\begin{itemize}
  \item \texttt{public\_standards\_to\_m\&a\_claims.md}: Board-level assertion table;
        Slide-to-Receipt bridge formula; buyer replication requirements.
  \item \texttt{public\_standards\_to\_blue\_river\_dam.md}: Validation gate execution
        protocol; dynamic threshold enforcement; POWL Projection Law.
  \item \texttt{STANDARDS\_TO\_BOARD\_CLAIMS.md}: Cross-walk from standards metrics
        to board-admissible assertions.
  \item \texttt{public\_standards\_to\_lifecycle\_actuation.md}: Actuation transaction
        JSON with \texttt{preflight\_soundness\_proof} signature.
\end{itemize}

\paragraph{Adversarial and verification files.}
\begin{itemize}
  \item \texttt{ocel-adversarial-gravity.md}: Adversarial test report OCEL-XES-ADV-V30.1.1.
  \item \texttt{xes-cryptographic-tombstoning.md}: XES GDPR/RTBF compliance gap analysis;
        \texttt{Event<Encrypted<PII>> $\to$ Event<Tombstoned<BLAKE3\_Receipt>>} typestate
        injection specification.
  \item \texttt{wf-net\_verification\_specification.md} v30.1.2: Karp-Miller Coverability
        Graph construction; Short-Circuited Petri Net analysis; Rust engine blueprint.
\end{itemize}

\subsection{Commands}
\label{sec:standards-impl-commands}

The standards project does not expose a build artifact or CLI binary. Its outputs are
markdown placement files, JSON schemas, and checkpoint ledger entries. The canonical
verification command is to run the conformance checker in wasm4pm against an event log
with a placement file's thresholds as configuration. The audit matrix at
\texttt{audit\_\_standards\_coverage.md} serves as the human-readable verification record;
the checkpoint JSON block in \texttt{checkpoint\_\_public\_standards\_crosswalk\_complete.md}
serves as the machine-readable ledger entry.

\subsection{Data and Ontology Surfaces}
\label{sec:standards-impl-data}

\paragraph{JSON Schema surface.}
The OTel trace integration schema (\texttt{otel\_trace\_integration\_schema.json}) is a
machine-validatable JSON Schema. It defines the \texttt{blake3\_receipt} field as a
required property of each OTel span, typed as a 64-character hexadecimal string. The
\texttt{event\_chain\_root} property captures the BLAKE3 hash of the entire trace, providing
a single-field integrity check for the complete span sequence.

\paragraph{Ledger block schema.}
The checkpoint file defines the ledger block schema with the following required fields:
\texttt{checkpoint\_id}, \texttt{timestamp}, \texttt{status}, \texttt{standards\_count\_verified},
\texttt{standards\_directory\_hash} (BLAKE3), and \texttt{witness\_signature} (ED25519).
The WF-net decommissioning receipt example extends this schema with process-specific
fields: \texttt{historical\_replay\_fitness: 0.978}, \texttt{precision: 0.912}, and
\texttt{total\_recorded\_cases: 124500}.

\paragraph{LossReport schema.}
Every trans-standard conversion emits a JSON \texttt{LossReport} with: \texttt{loss\_report\_id},
\texttt{timestamp}, \texttt{source\_format}, \texttt{target\_format}, \texttt{structural\_changes}
(sub-object naming per-conversion-type quantities such as discarded O2O links, pruned
E2O links, synthetic objects materialized), and \texttt{witness\_signature}.

\subsection{Templates and Rendered Artifacts}
\label{sec:standards-impl-templates}

The project's primary rendered artifacts are the placement files themselves---each is a
fully realized document rather than a template. The \texttt{docs-law\_\_standards\_readme.md}
defines the ledger block schema and serves as the authoritative reference for the
directory's structure and conventions.

The ggen projection map (\texttt{public\_standards\_to\_ggen\_projections.md}) defines
the mapping from standard constructs to Rust code artifacts: Petri net places project
to Rust structs, BPMN gateways project to \texttt{match}/\texttt{join!} patterns, and
Declare templates project to FSA implementations. These projections are the bridge from
the standards registry to executable wasm4pm code.
